\documentclass[11pt]{article}
\usepackage[margin=0.82in]{geometry}
\usepackage[T1]{fontenc}
\usepackage{amsmath,amssymb,amsthm,mathtools}
\usepackage{graphicx}
\usepackage{microtype}

\newtheorem{theorem}{Theorem}
\newcommand{\R}{\mathbb{R}}
\newcommand{\M}{\mathcal{M}}
\newcommand{\Hp}{\mathcal{H}^{+}}

\title{Problem A is the finite Bessel-energy class\\
\large A literature-implied full characterization for arXiv:1412.5855}
\author{Run \texttt{fa\_banach\_001} --- \texttt{agent\_lane\_03}}
\date{13 August 2026}

\begin{document}
\maketitle

\noindent\textbf{Status.}
Full answer to the stated characterization problem, but not an originality
claim.  It is a direct specialization of classical Bessel-potential theory.

\section*{The source problem}

Jamr\'oz writes $\M_+(\R^2)$ for the finite nonnegative Radon measures and
defines $H^{-1}(\R^2)$ by the Fourier weight $(1+|\xi|^2)^{-1}$.  Problem A
asks for a characterization of their intersection.  The source then quotes a
positive-logarithmic-energy criterion under a compact-support hypothesis.

\begin{center}
  \includegraphics[width=0.82\textwidth]{figures/source_problem_crop.png}
\end{center}

The exact source pages are bundled as
\texttt{source\_excerpt\_problem\_A\_pages\_1\_2.pdf}.

For a finite nonnegative measure $\mu$, put
\[
 \Hp(\mu)=\iint_{\R^2\times\R^2}
       \log^+\!\frac1{|x-y|}\,d\mu(x)d\mu(y).
\]

\begin{theorem}[Global characterization]
For every finite nonnegative Radon measure $\mu$ on $\R^2$,
\[
 \mu\in H^{-1}(\R^2) \quad\Longleftrightarrow\quad \Hp(\mu)<\infty.
\]
Thus the same criterion characterizes all of
$\M_+(\R^2)\cap H^{-1}(\R^2)$; compact support is unnecessary.
\end{theorem}

\section*{Proof}

Use the Fourier convention
\[
 \widehat f(\xi)=\int_{\R^2}e^{-ix\cdot\xi}f(x)\,dx,
 \qquad
 f(x)=(2\pi)^{-2}\int_{\R^2}e^{ix\cdot\xi}\widehat f(\xi)\,d\xi.
\]
Let
\[
 h_t(x)=\frac1{4\pi t}e^{-|x|^2/(4t)},\qquad
 G(x)=\int_0^\infty e^{-t}h_t(x)\,dt.
\]
Then $\widehat h_t(\xi)=e^{-t|\xi|^2}$ and hence
\[
 \widehat G(\xi)=\int_0^\infty e^{-t(1+|\xi|^2)}\,dt
 =\frac1{1+|\xi|^2}.
\]
For fixed $t>0$, Fourier inversion for the bounded heat kernel gives
\[
 \iint h_t(x-y)\,d\mu(x)d\mu(y)
 =(2\pi)^{-2}\int_{\R^2}e^{-t|\xi|^2}|\widehat\mu(\xi)|^2\,d\xi.
\]
All quantities are nonnegative.  Tonelli's theorem therefore yields the
extended-real identity
\begin{equation}
 \iint G(x-y)\,d\mu(x)d\mu(y)
 =(2\pi)^{-2}\int_{\R^2}
     \frac{|\widehat\mu(\xi)|^2}{1+|\xi|^2}\,d\xi.       \tag{1}
\end{equation}
In particular, $\mu\in H^{-1}(\R^2)$ exactly when its $G$-energy is finite.

The heat integral is the order-two Bessel kernel in dimension two:
\[
 G(x)=\frac1{2\pi}K_0(|x|),
\]
where $K_0$ is the modified Bessel function.  The classical asymptotics are
\[
 K_0(r)\sim\log\frac1r\quad(r\downarrow0),
 \qquad
 K_0(r)\sim\sqrt{\frac{\pi}{2r}}e^{-r}\quad(r\to\infty).
\]
Consequently there is a constant $C$ such that, for every finite positive
$\mu$ of mass $m$,
\[
 \Hp(\mu)\le C\iint G(x-y)\,d\mu(x)d\mu(y),
\]
and
\[
 \iint G(x-y)\,d\mu(x)d\mu(y)
 \le C\bigl(\Hp(\mu)+m^2\bigr).
\]
Indeed, on $0<r<1$ the two singular kernels are comparable, while on
$r\ge1$ the Bessel kernel is bounded.  Finiteness of the two energies is
therefore equivalent.  Combining this with (1) proves the theorem. \qed

\section*{Proof Intuition}

The inhomogeneous Sobolev weight is already a positive-definite interaction
kernel.  Writing its reciprocal as a Laplace transform of heat multipliers
turns the Fourier norm into a manifestly positive pair energy.  In two
dimensions that interaction is $K_0$: logarithmic at short range and rapidly
decaying at long range.  A finite measure has no difficulty at long range, so
only the positive logarithmic self-interaction decides membership.

\section*{Literature relation, scope, and review recommendation}

Aronszajn and Smith give the Bessel kernel's Fourier multiplier and explicit
modified-Bessel formula on printed pp.~413--414, and the $K_0$ asymptotics on
printed p.~415.  Those pages are bundled in
\texttt{supporting\_excerpt\_bessel\_kernel\_printed\_413\_415.pdf}.
The theorem above is therefore a direct two-dimensional finite-measure
specialization of their framework, not a priority claim.

The conclusion fully answers Problem A as phrased.  It does not solve the
separate Euler uniqueness problem mentioned as motivation.  Human review can
focus on the extended-real identity (1), the source's meaning of ``bounded
Radon measure'' as finite mass, and the two kernel comparisons.

\section*{References}

\begin{enumerate}
\item G. Jamr\'oz, \emph{Nonnegative measures belonging to
  $H^{-1}(\R^2)$}, arXiv:1412.5855, Problem A on PDF p.~1 and Lemma 1.1 on
  PDF p.~2.
\item N. Aronszajn and K. T. Smith, \emph{Theory of Bessel Potentials. I},
  Ann. Inst. Fourier \textbf{11} (1961), 385--475; formulas (2.7)--(2.10),
  (3.4), and (3.5), printed pp.~414--415.
\end{enumerate}

\end{document}
